\documentclass[11pt]{article}
\usepackage[a4paper, top=18mm, bottom=18mm, left=20mm, right=20mm]{geometry}
\usepackage[T2A]{fontenc}
\usepackage[utf8]{inputenc}
\usepackage[ukrainian]{babel}
\usepackage{graphicx}
\usepackage[export]{adjustbox}
\usepackage{amsmath}
\usepackage{amssymb}
\graphicspath{{/app/Quiz/Scripts/2023/ProgAll/15BinTree/}{/app/Quiz/Scripts/2021/Mod2021_3/BinTree/}{/app/Quiz/Scripts/2020/Mod2020_4/BinTree/}}
\setlength{\parindent}{0pt}
\setlength{\parskip}{6pt}
\pagestyle{empty}
\begin{document}
%begin
{\large\textbf{Контрольна робота}}

\textbf{Варіант \No\,1}\hfill \textit{ПІБ:}\,\underline{\hspace{60mm}}
\vspace{4mm}

%beginitem
1. %goodpath:Scripts/2018/Mod2018_1/03-good.txt
Відмітити все, що є коректним оператором С++:
( 1 ) return x<=2;

( 2 ) double n.m;

( 3 ) cin>>x+3;

( 4 ) if (1>2) x=10;
\vspace{5mm}
%enditem
%beginitem
2. Що буде виведено за виконання фрагменту коду?

%code
3**-1--1
%code
\vspace{5mm}
%enditem
%beginitem
3. Що буде виведено за виконання фрагменту коду?

%code
print(3.0 == 3 is 6 != 9)
%code
\vspace{5mm}
%enditem
%beginitem
4. %goodpath:Scripts/2019/Mod2019_1/Lex/01-good.txt
Відмітити все, що є однією правильною лексемою:
( 1 ) 'c'

( 2 ) 0123

( 3 ) 0b23

( 4 ) xy
\vspace{5mm}
%enditem
%beginitem
5. 

a,b,c=1,4,7
def g(b):
global c    a-=5
    b=1
    c=3
    return a+b+c

a,b,c=0,9,5
print(g(a),a,b,c)
\vspace{5mm}
%enditem
%beginitem
6. Що буде виведено за виконання фрагменту коду?

%code
cout << (8 / 8 / 8 / 8);
%code
\vspace{5mm}
%enditem
%beginitem
7. 

\begin{center}
	\adjustimage{max size={0.8\linewidth}{0.8\paperheight}}
	{../15BinTree/trees_exp/35.png}
\end{center}
Указати послідовність міток вершин дерева, зображеного на рис., що утвориться в результаті обходу дерева в оберненому порядку. Мітки записувати через пробіл.\par Алгоритм починає роботу з кореня (на рис. це верхня вершина).
%answer:35 оберненому
\vspace{5mm}
%enditem
%beginitem
8. %goodpath:Scripts/2019/Mod2019_1/Lex/03-good.txt
Відмітити все, що є коректним оператором С++:
( 1 ) return 'a'<'c'

( 2 ) print(sep=":","qwe")

( 3 ) x=*2

( 4 ) float(2)/87
\vspace{5mm}
%enditem
%beginitem
9. Що буде виведено за виконання фрагменту коду?
%code
a, b, c = 7, 9, 8
def g(a, b, c=6):
    print(a, b, c, end=" ")

g(1, 4, 2)
print(a, b, c)
%code
\vspace{5mm}
%enditem
%beginitem
10. 

print('R', end=' ')
b=['0','1','2','3','4']
b.append('13')
print(b)

\vspace{5mm}
%enditem










%end
\newpage

\newpage

%begin
{\large\textbf{Контрольна робота}}

\textbf{Варіант \No\,2}\hfill \textit{ПІБ:}\,\underline{\hspace{60mm}}
\vspace{4mm}

%beginitem
1. 

a,b,c=1,4,7
def g(b):
global c    a-=5
    b=1
    c=3
    return a+b+c

a,b,c=0,9,5
print(g(a),a,b,c)
\vspace{5mm}
%enditem
%beginitem
2. Що буде виведено за виконання фрагменту коду?

%code
print('R', end=' ')
b = ('a','b','c','d','e',0,1,2,3,4)
print(b[-1])
%code

\vspace{5mm}
%enditem
%beginitem
3. %goodpath:Scripts/2018/Mod2018_1/03-good.txt
Відмітити все, що є коректним оператором С++:
( 1 ) return x<=2;

( 2 ) double n.m;

( 3 ) cin>>x+3;

( 4 ) if (1>2) x=10;
\vspace{5mm}
%enditem
%beginitem
4. Що буде виведено за виконання фрагменту коду?

%code
cout << (!3.0<=3 or !6== 9 or 6.0>false);
%code
\vspace{5mm}
%enditem
%beginitem
5. Що буде виведено за виконання фрагменту коду?
%code

a,b,c=7,9,8
def g(a,b,c=6):
    print(a,b,c,end="")

g(1,4,2)
print(a,b,c)

%code
\vspace{5mm}
%enditem
%beginitem
6. 

x,y,z=1,4,2
y,y,z,x=y,x,7,z
print(x,y,z)
\vspace{5mm}
%enditem
%beginitem
7. %goodpath:Scripts/2019/Mod2019_1/Lex/03-good.txt
Відмітити все, що є коректним оператором С++:
( 1 ) return 'a'<'c'

( 2 ) print(sep=":","qwe")

( 3 ) x=*2

( 4 ) float(2)/87
\vspace{5mm}
%enditem
%beginitem
8. %goodpath:Scripts/2018/Mod2018_1/01-good.txt
Відмітити все, що є однією правильною лексемою:
( 1 ) a2f

( 2 ) 'qwerty'

( 3 ) **

( 4 ) true
\vspace{5mm}
%enditem
%beginitem
9. Що буде виведено за виконання фрагменту коду?

%code
print(3**-1--1)
%code
\vspace{5mm}
%enditem
%beginitem
10. Що буде виведено за виконання фрагменту коду?

%code
#include <iostream>
using namespace std;

int g(int b){
    int v = 45;
    if (b) 
        v = 1;
    else if (b <= -1)
         return 4;
    else
         v = 7;
    return v;
}

int main(){
    cout << g(3);
    return 0;
}
%code

\vspace{5mm}
%enditem










%end
\newpage

\newpage

%begin
{\large\textbf{Контрольна робота}}

\textbf{Варіант \No\,3}\hfill \textit{ПІБ:}\,\underline{\hspace{60mm}}
\vspace{4mm}

%beginitem
1. Що буде виведено за виконання фрагменту коду?

%code
cout << 8 / 8 / 8 / 8;
%code
\vspace{5mm}
%enditem
%beginitem
2. %goodpath:Scripts/2019/Mod2019_1/Lex/01-good.txt
Відмітити все, що є однією правильною лексемою:
( 1 ) 'c'

( 2 ) 0123

( 3 ) 0b23

( 4 ) xy
\vspace{5mm}
%enditem
%beginitem
3. 
try:
    res = 3**-1--1
    if isinstance(res, bool):
        print(f'bool;{res}')
    elif isinstance(res, int):
        print(f'int;{res}')
    elif isinstance(res, float):
        print(f'float;{res:.2f}')
    else:
        print('???')
except:
    print('Z')
\vspace{5mm}
%enditem
%beginitem
4. 
try:
    res = 7//7/7*7
    if isinstance(res, bool):
        print(f'bool;{res}')
    elif isinstance(res, int):
        print(f'int;{res}')
    elif isinstance(res, float):
        print(f'float;{res:.2f}')
    else:
        print('???')
except:
    print('Z')
\vspace{5mm}
%enditem
%beginitem
5. 
5.0<="9.0j"\vspace{5mm}
%enditem
%beginitem
6. %goodpath:Scripts/2019/Mod2019_1/Lex/03-good.txt
Відмітити все, що є коректним оператором С++:
( 1 ) return 'a'<'c'

( 2 ) print(sep=":","qwe")

( 3 ) x=*2

( 4 ) float(2)/87
\vspace{5mm}
%enditem
%beginitem
7. 
try:
    for b in range(-3, -3, 3):
        if b<=4:
            continue
        print(b, end=';');b = 2
        if b<=4:
            break
    else:
        print(b,end=';')
    print(b)
except:
    print('Z')
\vspace{5mm}
%enditem
%beginitem
8. Що буде виведено за виконання фрагменту коду?
%code
#include <iostream>
using namespace std;

int f(int &x, int &y){
    x = 5;
    y-= 9;
    return y;
}

int main(){
    int a = 6, b = 4;
    b = f(b, b);
    cout << a << ":" << b <<':';
    {
        int a = 2, b = 3;
        cout << ((b<=7) || ((a-=2) <= 7))
             << a << ":" << b << ":";
    }
    {
        inta = 1;
        cout << a << ':';
    }
    cout << a << ":" << b << endl;
    return 0;
}
%code

\vspace{5mm}
%enditem
%beginitem
9. %goodpath:Scripts/2018/Mod2018_1/01-good.txt
Відмітити все, що є однією правильною лексемою:
( 1 ) a2f

( 2 ) 'qwerty'

( 3 ) **

( 4 ) true
\vspace{5mm}
%enditem
%beginitem
10. Що буде виведено за виконання фрагменту коду?

%code
print('R', end=' ')
j=1
while j<=9:
    j+=2
print(j)
%code

\vspace{5mm}
%enditem










%end
\newpage

\newpage

%begin
{\large\textbf{Контрольна робота}}

\textbf{Варіант \No\,4}\hfill \textit{ПІБ:}\,\underline{\hspace{60mm}}
\vspace{4mm}

%beginitem
1. 
a = 1
b = 4
c = 7

def g(b):
global c    a -= 5
    b = 1
    c = 3
    return a + b + c

a = 0
b = 9
c = 5

try:
    print(g(a), a, b, c, sep=';')
except:
    print('Z', a, b, c, sep=';')
\vspace{5mm}
%enditem
%beginitem
2. 
Записати ВИРАЗ, значенням якого є множина, що складається з цілих чисел від 14 до 2012 (включно), що діляться на 3.\vspace{5mm}
%enditem
%beginitem
3. %goodpath:Scripts/2019/Mod2019_1/Lex/03-good.txt
Відмітити все, що є коректним оператором С++:
( 1 ) return 'a'<'c'

( 2 ) print(sep=":","qwe")

( 3 ) x=*2

( 4 ) float(2)/87
\vspace{5mm}
%enditem
%beginitem
4. 
Записати ВИРАЗ, значенням якого є список, що складається з квадратних коренівцілих чисел від 14 до 2012 (включно), що діляться на 3, записаних в оберненому порядку.\vspace{5mm}
%enditem
%beginitem
5. %goodpath:Scripts/2018/Mod2018_1/01-good.txt
Відмітити все, що є однією правильною лексемою:
( 1 ) a2f

( 2 ) 'qwerty'

( 3 ) **

( 4 ) true
\vspace{5mm}
%enditem
%beginitem
6. 
Написати програму, що запитує у користувача значення дійсних параметрів $x$ та $y$, обчислює для них значення виразу 
$$\frac{\pi x - 19}{(\arcsin x - 0.5) (y - 23)}$$ 
та повiдомляє введенi данi та результат обчислення.
 Оцінюється правильність та якість коду.

\vspace{5mm}
%enditem
%beginitem
7. %goodpath:Scripts/2019/Mod2019_1/Lex/01-good.txt
Відмітити все, що є однією правильною лексемою:
( 1 ) 'c'

( 2 ) 0123

( 3 ) 0b23

( 4 ) xy
\vspace{5mm}
%enditem
%beginitem
8. 

print('R', end=' ')
for b in range(-3,-3,3):
    if b<=4:
        continue
        print(b, end=' ')
        b=2
    if b<=4:
        break
    else:
        print(b, end=' ')
    print(b)
\vspace{5mm}
%enditem
%beginitem
9. 
b=['0','1','2','3','4']; b.append('13'); print(b)\vspace{5mm}
%enditem
%beginitem
10. 

\begin{center}
	\adjustimage{max size={0.8\linewidth}{0.8\paperheight}}
	{../15BinTree/trees_exp/35.png}
\end{center}
Указати послідовність міток вершин дерева, зображеного на рис., що утвориться в результаті обходу дерева в оберненому порядку. Мітки записувати через пробіл.\par Алгоритм починає роботу з кореня (на рис. це верхня вершина).
%answer:35 оберненому
\vspace{5mm}
%enditem










%end
\newpage

\newpage

%begin
{\large\textbf{Контрольна робота}}

\textbf{Варіант \No\,5}\hfill \textit{ПІБ:}\,\underline{\hspace{60mm}}
\vspace{4mm}

%beginitem
1. Що буде виведено за виконання фрагменту коду?

%code
print(3**-1--1)
%code
\vspace{5mm}
%enditem
%beginitem
2. %goodpath:Scripts/2019/Mod2019_1/Lex/01-good.txt
Відмітити все, що є однією правильною лексемою:
( 1 ) 'c'

( 2 ) 0123

( 3 ) 0b23

( 4 ) xy
\vspace{5mm}
%enditem
%beginitem
3. 

\begin{center}
	\adjustimage{max size={0.8\linewidth}{0.8\paperheight}}
	{../BinTree/trees_exp/35.png}
\end{center}
Указати послідовність міток вершин дерева, зображеного на рис., що утвориться в результаті обходу дерева в оберненому порядку. Мітки записувати через пробіл.\par Алгоритм починає роботу з кореня (на рис. це верхня вершина).
%answer:35 оберненому
\vspace{5mm}
%enditem
%beginitem
4. 

print('R', end=' ')
j=1
while j<=9:
    j+=2
print(j)

\vspace{5mm}
%enditem
%beginitem
5. 
Записати ВИРАЗ, значенням якого є множина, що складається з квадратних коренівцілих чисел від 14 до 2012 (включно), що діляться на 3.\vspace{5mm}
%enditem
%beginitem
6. %goodpath:Scripts/2019/Mod2019_1/Lex/03-good.txt
Відмітити все, що є коректним оператором С++:
( 1 ) return 'a'<'c'

( 2 ) print(sep=":","qwe")

( 3 ) x=*2

( 4 ) float(2)/87
\vspace{5mm}
%enditem
%beginitem
7. 
try:
    j = 1
    while j<=9:
        j += 2
    print(j, end=';')
except:
    print('ERR')
\vspace{5mm}
%enditem
%beginitem
8. %goodpath:Scripts/2018/Mod2018_1/01-good.txt
Відмітити все, що є однією правильною лексемою:
( 1 ) a2f

( 2 ) 'qwerty'

( 3 ) **

( 4 ) true
\vspace{5mm}
%enditem
%beginitem
9. Що буде виведено за виконання фрагменту коду?

%code
#include <iostream>
using namespace std;

int a = 1, b = 4, c = 7;

int g(){
    inta = 1;
    intb = 0;
    intc = 9;
    return a + b + c;
}

int main(){
    inta = 5;
    intb = 6;
    intc = 8;
    cout << g() << ':';
    cout << a << ':' << b << ':' << c;
    return 0;
}
%code
\vspace{5mm}
%enditem
%beginitem
10. Що буде виведено за виконання фрагменту коду?

%code
a,b,c=1,4,7
def g(b):
global c    a-=5
    b=1
    c=3
    return a+b+c

a,b,c=0,9,5
print(g(a),a,b,c)
%code
\vspace{5mm}
%enditem










%end
\newpage

\newpage

%begin
{\large\textbf{Контрольна робота}}

\textbf{Варіант \No\,6}\hfill \textit{ПІБ:}\,\underline{\hspace{60mm}}
\vspace{4mm}

%beginitem
1. Що буде виведено за виконання фрагменту коду?

%code
a,b,c=1,4,7
def g(b):
global c    a-=5
    b=1
    c=3
    return a+b+c

a,b,c=0,9,5
print(g(a),a,b,c)
%code
\vspace{5mm}
%enditem
%beginitem
2. %goodpath:Scripts/2019/Mod2019_1/Lex/03-good.txt
Відмітити все, що є коректним оператором С++:
( 1 ) return 'a'<'c'

( 2 ) print(sep=":","qwe")

( 3 ) x=*2

( 4 ) float(2)/87
\vspace{5mm}
%enditem
%beginitem
3. Що буде виведено за виконання фрагменту коду?

%code
for b in range(2,2+6,-2):
    if b<=4:
        continue
    print(b, end=' ')
    b=2
    if b<=4:
        break
else:
    print(b, end=' ')
print(b, end=' ')
%code
\vspace{5mm}
%enditem
%beginitem
4. Що буде виведено за виконання фрагменту коду?

%code
print(3.0==3==6==9)
%code
\vspace{5mm}
%enditem
%beginitem
5. Що буде виведено за виконання фрагменту коду?

%code
a = 1
b = 4
c = 7

def g(b):
global c    a=5
    b=1
    c=3
    return a+b+c

a, b, c = 0,9,5
try:
    print(g(a), a, b, c, sep=';')
except:
    print('Z', a, b, c, sep=';')
%code
\vspace{5mm}
%enditem
%beginitem
6. Що буде виведено за виконання фрагменту коду?

%code
print('R', end=' ')
j=9
while j>=1:
    j+=-2
print(j)
%code

\vspace{5mm}
%enditem
%beginitem
7. Що буде виведено за виконання фрагменту коду?

%code
def g(b):
    v=45
    if b: 
        v=1
    elif b<=-1:
         return 4
    else:
         v=7
    return v

print(g(3))
%code
\vspace{5mm}
%enditem
%beginitem
8. Що буде виведено за виконання фрагменту коду?

%code
not3.0<=3 or not6==9 or 6.0>False
%code
\vspace{5mm}
%enditem
%beginitem
9. %goodpath:Scripts/2018/Mod2018_1/03-good.txt
Відмітити все, що є коректним оператором С++:
( 1 ) return x<=2;

( 2 ) double n.m;

( 3 ) cin>>x+3;

( 4 ) if (1>2) x=10;
\vspace{5mm}
%enditem
%beginitem
10. %goodpath:Scripts/2019/Mod2019_1/Lex/01-good.txt
Відмітити все, що є однією правильною лексемою:
( 1 ) 'c'

( 2 ) 0123

( 3 ) 0b23

( 4 ) xy
\vspace{5mm}
%enditem










%end
\newpage

\newpage

%begin
{\large\textbf{Контрольна робота}}

\textbf{Варіант \No\,7}\hfill \textit{ПІБ:}\,\underline{\hspace{60mm}}
\vspace{4mm}

%beginitem
1. 

def g(b):
    v=45
    if b: 
        v=1
    elif b<=-1:
         return 4
    else:
         v=7
    return v

print(g(3))
\vspace{5mm}
%enditem
%beginitem
2. %goodpath:Scripts/2019/Mod2019_1/Lex/03-good.txt
Відмітити все, що є коректним оператором С++:
( 1 ) return 'a'<'c'

( 2 ) print(sep=":","qwe")

( 3 ) x=*2

( 4 ) float(2)/87
\vspace{5mm}
%enditem
%beginitem
3. Що буде виведено за виконання фрагменту коду?

%code
a,b,c=1,4,7
def g(b):
    a=5
    b=1
    c=3
    return a+b+c

a,b,c=0,9,5
print(g(a),a,b,c)
%code
\vspace{5mm}
%enditem
%beginitem
4. 
Написати програму, що запитує у користувача значення дійсних параметрів $x$ та $y$, обчислює для них значення виразу та повiдомляє введенi данi та результат обчислення.
$$\frac{\pi x - 19}{(\arcsin x - 0.5) (y - 23)}$$ 

\vspace{5mm}
%enditem
%beginitem
5. %goodpath:Scripts/2018/Mod2018_1/01-good.txt
Відмітити все, що є однією правильною лексемою:
( 1 ) a2f

( 2 ) 'qwerty'

( 3 ) **

( 4 ) true
\vspace{5mm}
%enditem
%beginitem
6. Що буде виведено за виконання фрагменту коду?

%code
j=1
while j<=9:
    j+=2
print(j)
%code

\vspace{5mm}
%enditem
%beginitem
7. Що буде виведено за виконання фрагменту коду?

%code
def g(a,b):
    c=45
    if a:
        c=1
    elif b<=-1:
         return 4
    else: 
        c=7
    return c

print(g(3,3))
%code
\vspace{5mm}
%enditem
%beginitem
8. %goodpath:Scripts/2018/Mod2018_1/03-good.txt
Відмітити все, що є коректним оператором С++:
( 1 ) return x<=2;

( 2 ) double n.m;

( 3 ) cin>>x+3;

( 4 ) if (1>2) x=10;
\vspace{5mm}
%enditem
%beginitem
9. 

def f(n):
    print("f", end="");
    return n<=0

print(f(3) or f(-1))
\vspace{5mm}
%enditem
%beginitem
10. Що буде виведено за виконання фрагменту коду?

%code
cout << (7 / 7 / 7 / 7);
%code
\vspace{5mm}
%enditem










%end
\newpage

\newpage

%begin
{\large\textbf{Контрольна робота}}

\textbf{Варіант \No\,8}\hfill \textit{ПІБ:}\,\underline{\hspace{60mm}}
\vspace{4mm}

%beginitem
1. Що буде виведено за виконання фрагменту коду?
%code
a, b, c = 7, 9, 8
def g(a, b, c=6):
    print(a, b, c, end=" ")

g(1, 4, 2)
print(a, b, c)
%code
\vspace{5mm}
%enditem
%beginitem
2. Що буде виведено за виконання фрагменту коду?

%code
j=1
while j<=9:
    j+=2
print(j)
%code

\vspace{5mm}
%enditem
%beginitem
3. 

def f(n):
    print("f", end="");
    return n<=0

print(f(3) or f(-1))
\vspace{5mm}
%enditem
%beginitem
4. %goodpath:Scripts/2019/Mod2019_1/Lex/01-good.txt
Відмітити все, що є однією правильною лексемою:
( 1 ) 'c'

( 2 ) 0123

( 3 ) 0b23

( 4 ) xy
\vspace{5mm}
%enditem
%beginitem
5. Що буде виведено за виконання фрагменту коду?
%code
#include <iostream>
using namespace std;

int f(int &x, int &y){
    x = 5;
    y-= 9;
    return y;
}

int main(){
    int a = 6, b = 4;
    b = f(b, b);
    cout << a << ":" << b <<':';
    {
        int a = 2, b = 3;
        cout << ((b<=7) || ((a-=2) <= 7)) << ':';
        cout << a << ":" << b << ':';
    }
    {
        inta = 1;
        cout << a << ':';
    }
    cout << a << ":" << b << endl;
    return 0;
}
%code

\vspace{5mm}
%enditem
%beginitem
6. Що буде виведено за виконання фрагменту коду?

%code
print(7//7/7*7)
%code
\vspace{5mm}
%enditem
%beginitem
7. %goodpath:Scripts/2019/Mod2019_1/Lex/03-good.txt
Відмітити все, що є коректним оператором С++:
( 1 ) return 'a'<'c'

( 2 ) print(sep=":","qwe")

( 3 ) x=*2

( 4 ) float(2)/87
\vspace{5mm}
%enditem
%beginitem
8. Що буде виведено за виконання фрагменту коду?

%code
print('R', end=' ')
b=['0','1','2','3','4']
b[-8:-8]='12'
print(b)
%code

\vspace{5mm}
%enditem
%beginitem
9. 
not3.0<=3 or not6==9 or 6.0>False\vspace{5mm}
%enditem
%beginitem
10. %goodpath:Scripts/2018/Mod2018_1/01-good.txt
Відмітити все, що є однією правильною лексемою:
( 1 ) a2f

( 2 ) 'qwerty'

( 3 ) **

( 4 ) true
\vspace{5mm}
%enditem










%end
\newpage

\newpage

%begin
{\large\textbf{Контрольна робота}}

\textbf{Варіант \No\,9}\hfill \textit{ПІБ:}\,\underline{\hspace{60mm}}
\vspace{4mm}

%beginitem
1. Що буде виведено за виконання фрагменту коду?
%code
if (19 <= 19)
    cout << "t";
else
    cout << "e";
%code

\vspace{5mm}
%enditem
%beginitem
2. Що буде виведено за виконання фрагменту коду?
code
__b = 0

class B:
    __b = 1
    
    def __init__(self):
        self.__b = 2
        __b = 3
        B.__b = 4

obj = B()
try:
    print(__b, end=' ')
    print(obj.__b, end=' ')
    print(B.__b, end=' ')
    print(obj.__b, end=' ')
    print(obj._B__b, end=' ')
except:
    print('ERR')
code
\vspace{5mm}
%enditem
%beginitem
3. %goodpath:Scripts/2019/Mod2019_1/Lex/03-good.txt
Відмітити все, що є коректним оператором С++:
( 1 ) return 'a'<'c'

( 2 ) print(sep=":","qwe")

( 3 ) x=*2

( 4 ) float(2)/87
\vspace{5mm}
%enditem
%beginitem
4. %goodpath:Scripts/2018/Mod2018_1/03-good.txt
Відмітити все, що є коректним оператором С++:
( 1 ) return x<=2;

( 2 ) double n.m;

( 3 ) cin>>x+3;

( 4 ) if (1>2) x=10;
\vspace{5mm}
%enditem
%beginitem
5. Що буде виведено за виконання фрагменту коду?

%code
print('R', end=' ')
b=['0','1','2','3','4']
b[-8:-8]='12'
print(b)
%code

\vspace{5mm}
%enditem
%beginitem
6. Що буде виведено за виконання фрагменту коду?

%code
def g(a,b):
    c=45
    if a:
        c=1
    elif b<=-1:
         return 4
    else: 
        c=7
    return c

print(g(3,3))
%code
\vspace{5mm}
%enditem
%beginitem
7. Що буде виведено за виконання фрагменту коду?

%code
j=1
while j<=9:
    j+=2
print(j)
%code

\vspace{5mm}
%enditem
%beginitem
8. %goodpath:Scripts/2019/Mod2019_1/Lex/01-good.txt
Відмітити все, що є однією правильною лексемою:
( 1 ) 'c'

( 2 ) 0123

( 3 ) 0b23

( 4 ) xy
\vspace{5mm}
%enditem
%beginitem
9. Що буде виведено за виконання фрагменту коду?

%code
print('R', end=' ')
j=9
while j>=1:
    j+=-2
print(j)
%code

\vspace{5mm}
%enditem
%beginitem
10. Що буде виведено за виконання фрагменту коду?

%code
print('R', end=' ')
j=1
while j<=9:
    j+=2
print(j)
%code

\vspace{5mm}
%enditem










%end
\newpage

\newpage

%begin
{\large\textbf{Контрольна робота}}

\textbf{Варіант \No\,10}\hfill \textit{ПІБ:}\,\underline{\hspace{60mm}}
\vspace{4mm}

%beginitem
1. 
b=['0','1','2','3','4']; b[-8:-8]='12'; print(b)\vspace{5mm}
%enditem
%beginitem
2. 

\begin{center}
	\adjustimage{max size={0.8\linewidth}{0.8\paperheight}}
	{../BinTree/trees_exp/35.png}
\end{center}
Указати послідовність міток вершин дерева, зображеного на рис., що утвориться в результаті обходу дерева в оберненому порядку. Мітки записувати через пробіл.\par Обхід дерева починається з кореня (на рис. це верхня вершина).
%answer:35 оберненому
\vspace{5mm}
%enditem
%beginitem
3. %goodpath:Scripts/2018/Mod2018_1/01-good.txt
Відмітити все, що є однією правильною лексемою:
( 1 ) a2f

( 2 ) 'qwerty'

( 3 ) **

( 4 ) true
\vspace{5mm}
%enditem
%beginitem
4. 
try:
    print(1, end=';')
    print(int('4'), end=';')
    print(7, end=';')
except KeyboardInterrupt:
    print(0, end=';')
except Exception:
    print(9, end=';')
except:
    print('Z', end=';')
else:
    print(5, end=';')
finally:
    print(6)
\vspace{5mm}
%enditem
%beginitem
5. %goodpath:Scripts/2018/Mod2018_1/03-good.txt
Відмітити все, що є коректним оператором С++:
( 1 ) return x<=2;

( 2 ) double n.m;

( 3 ) cin>>x+3;

( 4 ) if (1>2) x=10;
\vspace{5mm}
%enditem
%beginitem
6. Що буде виведено за виконання фрагменту коду?

%code
print(7//7/7*7)
%code
\vspace{5mm}
%enditem
%beginitem
7. %goodpath:Scripts/2019/Mod2019_1/Lex/03-good.txt
Відмітити все, що є коректним оператором С++:
( 1 ) return 'a'<'c'

( 2 ) print(sep=":","qwe")

( 3 ) x=*2

( 4 ) float(2)/87
\vspace{5mm}
%enditem
%beginitem
8. Що буде виведено за виконання фрагменту коду?

%code
#include <iostream>

int g(int a, int b){
    int c = 45;
    if (a)
        c = 1;
    else if (b <= -1)
         return 4;
    else 
        c = 7;
    return c;
}

int main(){
    std::cout << g(3, 3);
    return 0;
}
%code
\vspace{5mm}
%enditem
%beginitem
9. 
try:
    res = 5.0j<="9.0"
    if isinstance(res, bool):
        print(f'bool;{res}')
    elif isinstance(res, int):
        print(f'int;{res}')
    elif isinstance(res, float):
        print(f'float;{res:.2f}')
    else:
        print('???')
except:
    print('Z')
\vspace{5mm}
%enditem
%beginitem
10. 
def f(n):
    print('f', end='')
    return n<=0

try:
    print(f(3) or f(-1))
except:
    print('ERR')
\vspace{5mm}
%enditem










%end
\newpage
\end{document}


